% !PRINTED 203
\weberpage{203}
%<<CONT
und wir erhalten eine zweite Möglichkeit:

\gesperrt{II. Diëdergruppe,} $\nu = \nu' = 2, \nu'' = m, n = 2m, m \geqq 2$

\[ \mu = \mu' = m, \quad \mu'' = 2. \]

Ist ferner keine der Zahlen $\nu', \nu''$ gleich 2, so muss eine von ihnen gleich 3 sein. Denn sind sie beide $\geqq 4$, so ist

\[ \mu' + \mu'' \leqq \frac{n}{2}, \]

was mit (16) im Widerspruche steht. Ist also $\nu' = 3$, $\mu' = \dfrac{n}{3}$, so folgt aus (16):

\weq{17}{\mu'' = \frac{n}{6} + 2, \quad \nu'' = \frac{6n}{n+12},}

woraus folgt, dass $\nu'' < 6$ sein muss, also nur einen der Werthe 3, 4, 5 haben kann.

Danach bekommen wir noch drei mögliche Fälle:

\gesperrt{III. Tetraëdergruppe:} $\nu = 2, \nu' = 3, \nu'' = 3, n = 12$

\[ \mu = 6, \quad \mu' = 4, \quad \mu'' = 4. \]

\gesperrt{IV. Octaëdergruppe:} $\nu = 2, \nu' = 3, \nu'' = 4, n = 24$

\[ \mu = 12, \quad \mu' = 8, \quad \mu'' = 6. \]

\gesperrt{V. Ikosaëdergruppe:} $\nu = 2, \nu' = 3, \nu'' = 5, n = 60$

\[ \mu = 30, \quad \mu' = 20, \quad \mu'' = 12. \]

Hiermit sind alle Möglichkeiten erschöpft.

Es ist freilich noch nicht bewiesen, dass diese Gruppen, die wir einstweilen mit den gebräuchlichen Namen aufgeführt haben, und die wir unter dem gemeinsamen Namen der \gesperrt{Polyëdergruppen} zusammenfassen, wirklich existiren, noch wie gross ihre Mannigfaltigkeit ist. Zu diesem Beweise führen erst die folgenden Betrachtungen.

\begin{center}
\S.\,53.

\gesperrt{Transformation der Substitutionen von $G$ auf einfache Formen.}
\end{center}

Ehe wir zur definitiven Aufstellung dieser Gruppen gehen, leiten wir einen Satz ab, der die Möglichkeit der vorkommenden Substitutionen noch weiter beschränkt.
